\section{The Library as a Test Oracle: Belenios}\label{sec:belenios}
The library has a second use that bears on the scheme. Instantiated at a
deployed system's group, encodings, and hash, its certified verifier is a
reference implementation of that system's proofs. Its certified prover
takes every random value and the witness as explicit inputs, so it can
produce proofs of \emph{false} statements with correctly derived
Fiat--Shamir challenges. These get past a verifier's hash checks and
exercise its algebraic checks. We applied this to
Belenios~\cite{belenios2019}, whose ballots carry the same disjunctive
encryption proofs as ours. We compared the audit tool its specification
prescribes, \texttt{belenios-tool election verify} (release 3.2.0), with
the certified verifier, on elections generated with the tool itself and on
1{,}715 honest, mutated, structural, and forged boards. The comparison
found two defects, both outside the proof verification.

First, the audit tool accepted 36 forged boards, each containing a single
invalid ballot inside an otherwise consistent archive: a well-formed
ciphertext of an out-of-range value, with random proofs and a signature
under a listed credential. The ballot was counted and the tool reported
``all checks passed''. One line of the tool forced the lazy
ballot-verification computation and discarded the resulting promise, so
the per-ballot proof and signature checks had no effect on the verdict.
Under the universal-verifiability threat model this lets a server stuff
or overwrite ballots while the official audit succeeds. The maintainers
confirmed the defect, present in releases 3.1 through 3.3.0, and
committed a fix (commit \texttt{115d0d47}, 22 October 2026, the day
after our report). The defect is in \texttt{belenios-tool} only and does
not affect the Belenios web server. Second, the shuffle verifier for
Belenios's mix-net questions, over Ed25519 (cofactor~8), applies the
subgroup test to the proof's commitments but not to its ciphertexts. A
mix server that adds a small-order point to one output and proves the
shuffle honestly is accepted for a challenge-dependent fraction of seeds
(26 of 100 for order~2), and the torsion survives the honest
re-encryptions of later servers \TODO{report this finding}. The
extracted certified shuffle verifier, with its own decoder's subgroup
check bypassed on purpose, accepts the same inputs at the same rate: the
verified algebra does not reject a torsion point, because in Rocq the
group type \emph{is} the prime-order subgroup and extraction erases that
proof. The certified pipeline rejects them all at its decoder.

Section~\ref{sec:discussion} applies the same lesson to our scheme: a
soundness theorem about a group type says nothing about deployed code
unless the decoder enforces the type, and deployed verifiers fail in the
checks around a correct proof verification.
